\Needspace{9\baselineskip}\begin{center}\textit{Teilungen des Kreises, die sich durch quadratische Gleichungen oder durch\\geometrische Konstruktionen ausführen lassen.}\end{center}


\Needspace{6\baselineskip}\Article{365}
Durch die vorstehenden Untersuchungen haben wir also die Teilung des Kreises in \(n\) Teile, wenn \(n\) eine Primzahl ist, auf die Lösung so vieler Gleichungen zurückgeführt, wie es Faktoren gibt, in die sich die Zahl \(n-1\) zerlegen lässt; die Grade dieser Gleichungen werden durch die Größe der Faktoren bestimmt. Ist daher \(n-1\) eine Potenz der Zahl 2, was für die Werte 3, 5, 17, 257, 65537, usw., von \(n\) eintritt, so führt die Teilung des Kreises auf lauter quadratische Gleichungen, und die trigonometrischen Funktionen der Winkel \(\frac Pn,\allowbreak \frac{2P}n\), usw., lassen sich durch mehr oder weniger komplizierte Quadratwurzeln (je nach der Größe von \(n\)) ausdrücken. In diesen Fällen lässt sich also die Teilung des Kreises in \(n\) Teile oder die Zeichnung eines regelmäßigen Vielecks mit \(n\) Seiten offenbar durch geometrische Kon\SourceLeaf{467}struktionen ausführen. So erhält man \(\text{z.\,B.}\) für \(n=17\) aus den Artikeln 354 und 361 leicht den folgenden Ausdruck für den Kosinus des Winkels \(\tfrac1{17}P\): \[\begin{aligned}&-\tfrac1{16}+\tfrac1{16}\surd17+\tfrac1{16}\surd(34-2\surd17)\\&\quad+\tfrac18\surd(17+3\surd17-\surd(34-2\surd17)\\&\qquad\qquad-2\surd(34+2\surd17))\end{aligned}\] Die Kosinus der Vielfachen dieses Winkels haben eine ähnliche Form, die Sinus aber ein Wurzelzeichen mehr. Es ist wahrlich sehr verwunderlich, dass, obwohl schon zu Euklids Zeiten die geometrische Teilbarkeit des Kreises in drei und fünf Teile bekannt war, diesen Entdeckungen während eines Zeitraums von 2000 Jahren nichts hinzugefügt wurde und alle Geometer es als gewiss verkündeten, dass sich außer diesen Teilungen und denjenigen, die sich von selbst daraus ergeben, nämlich den Teilungen in \(15,\allowbreak 3.2^\mu,\allowbreak 5.2^\mu,\allowbreak 15.2^\mu\) sowie in \(2^\mu\) Teile, keine weiteren durch geometrische Konstruktionen ausführen ließen. — Ferner lässt sich leicht beweisen, dass, wenn die Primzahl \(n\) \(=2^m+1\) ist, auch der Exponent \(m\) keine anderen Primfaktoren als die Zahl 2 enthalten kann und daher entweder \(=1\) oder \(=2\) oder einer höheren Potenz der Zahl 2 gleich sein muss. Wäre nämlich \(m\) durch irgendeine ungerade Zahl \(\zeta\) (größer als die Einheit) teilbar und \(m=\zeta\eta\), so wäre \(2^m+1\) durch \(2^\eta+1\) teilbar und daher notwendig zusammengesetzt. Alle Werte von \(n\), bei denen wir auf bloß quadratische Gleichungen geführt werden, sind also in der Form \(2^{2^\nu}+1\) enthalten; so erhält man die fünf Zahlen 3, 5, 17, 257, 65537, wenn man \(\nu=0,\allowbreak 1,\allowbreak 2,\allowbreak 3,\allowbreak 4\) oder \(m=1,\allowbreak 2,\allowbreak 4,\allowbreak 8,\allowbreak 16\) setzt. Keineswegs lässt sich aber die Teilung des Kreises für \textit{alle} in dieser Form enthaltenen Zahlen geometrisch ausführen, sondern nur für diejenigen, die Primzahlen sind. Zwar hatte Fermat, durch Induktion irregeführt, behauptet, alle in dieser Form enthaltenen Zahlen müssten Primzahlen sein; doch bemerkte der berühmte Euler als Erster, dass diese Regel schon für \(\nu=5\) oder \(m=32\) falsch ist, da die Zahl \(2^{32}+1=4294967297\) den Faktor 641 enthält.

Enthält dagegen \(n-1\) andere Primfaktoren als 2, so gelangen wir immer zu höheren Gleichungen: nämlich zu einer oder mehreren kubischen Gleichungen, wenn 3 einmal oder mehrmals unter den Primfaktoren von \(n-1\) vorkommt, zu Gleichungen fünften Grades, wenn \(n-1\) durch 5 teilbar ist, und so weiter; \textsc{und wir können mit voller Strenge beweisen, dass sich diese höheren Gleichungen in keiner Weise vermeiden oder auf niedrigere zurückführen lassen,} auch wenn uns die Grenzen dieses Werkes nicht erlauben, diesen Beweis hier zu geben. Wir glaubten dennoch, darauf hinweisen zu müssen, damit niemand hoffe, noch andere Teilungen als die von unserer Theorie nahegelegten, \(\text{z.\,B.}\) Teilungen in 7, 11, 13, 19, usw., Teile, auf geometrische Konstruktionen zurückzuführen und seine Zeit unnütz vergeude.
\SourceLeaf{468}


\Needspace{6\baselineskip}\Article{366}
Soll der Kreis in \(a^\alpha\) Teile geteilt werden, wobei \(a\) eine Primzahl bezeichnet, so lässt sich dies offenbar geometrisch ausführen, wenn \(a=2\), aber für keinen anderen Wert von \(a\), sofern \(\alpha>1\). Denn dann muss man außer den Gleichungen, die zur Teilung in \(a\) Teile erforderlich sind, notwendig noch \(\alpha-1\) weitere Gleichungen vom Grad \(a\) lösen; auch diese lassen sich in keiner Weise vermeiden oder im Grad senken. Die Grade der notwendigen Gleichungen lassen sich also allgemein aus den Primfaktoren der Zahl \((a-1)a^{\alpha-1}\) erkennen (also auch im Fall \(\alpha=1\)).

Soll schließlich der Kreis in \(N=a^\alpha b^\beta c^\gamma\ldots\) Teile geteilt werden, wobei \(a,\allowbreak b,\allowbreak c\), usw., verschiedene Primzahlen bezeichnen, so genügt es, die Teilungen in \(a^\alpha,\allowbreak b^\beta,\allowbreak c^\gamma\), usw., Teile ausgeführt zu haben (Artikel 336). Um daher die Grade der hierzu nötigen Gleichungen zu erkennen, muss man die Primfaktoren der Zahlen \[(a-1)a^{\alpha-1},\quad(b-1)b^{\beta-1},\quad(c-1)c^{\gamma-1}\ \text{usw.}\] oder, was hier auf dasselbe hinausläuft, ihres Produkts betrachten. Man beachte, dass dieses Produkt die Anzahl der Zahlen ausdrückt, die zu \(N\) teilerfremd und kleiner als diese Zahl sind (Artikel 38). Geometrisch lässt sich die Teilung also nur dann ausführen, wenn diese Anzahl eine Zweierpotenz ist; enthält sie dagegen andere Primfaktoren als \(2\), etwa \(p,\allowbreak p'\), usw., so lassen sich Gleichungen vom Grad \(p,\allowbreak p'\), usw., in keiner Weise vermeiden. Daraus folgt allgemein: Damit sich der Kreis geometrisch in \(N\) Teile teilen lässt, muss \(N\) \textit{entweder} \(2\) oder eine höhere Potenz von \(2\), \textit{oder} eine Primzahl der Form \(2^m+1\), \textit{oder} ein Produkt mehrerer solcher Primzahlen, \textit{oder} das Produkt aus einer oder mehreren solchen Primzahlen mit \(2\) oder einer höheren Potenz von \(2\) sein. Kürzer gesagt darf \(N\) weder einen ungeraden Primfaktor enthalten, der nicht die Form \(2^m+1\) hat, noch irgendeinen Primfaktor der Form \(2^m+1\) mehr als einmal. Solche Werte von \(N\) unterhalb von \(300\) sind die folgenden \(38\): \(2,\allowbreak 3,\allowbreak 4,\allowbreak 5,\allowbreak 6,\allowbreak 8,\allowbreak 10,\allowbreak 12,\allowbreak 15,\allowbreak 16,\allowbreak 17,\allowbreak 20,\allowbreak 24,\allowbreak 30,\allowbreak 32,\allowbreak 34,\allowbreak 40,\allowbreak 48,\allowbreak 51,\allowbreak 60,\allowbreak 64,\allowbreak 68,\allowbreak 80,\allowbreak 85,\allowbreak 96,\allowbreak 102,\allowbreak 120,\allowbreak 128,\allowbreak 136,\allowbreak 160,\allowbreak 170,\allowbreak 192,\allowbreak 204,\allowbreak 240,\allowbreak 255,\allowbreak 256,\allowbreak 257,\allowbreak 272.\)
\SourceLeaf{469}


